\documentclass[12pt]{article}
\usepackage[T1]{fontenc}
\usepackage[utf8]{inputenc}
\usepackage{lmodern}
\usepackage{textcomp}
\usepackage{enumitem}
\usepackage{geometry}
\usepackage{graphicx}
\usepackage{amsmath, amssymb}
\geometry{margin=1in}
\begin{document}
\begin{center}
Economics - Undergraduate Level Assessment 11 \
Total Marks: 40 \
Answer all questions.
\end{center}

\section*{Multiple Choice (10 marks)}
\begin{enumerate}[label=\arabic*.]
\item Which of the following is equal to one?
\begin{enumerate}[label=(\alph*)]
    \item The elasticity of the long-run aggregate supply curve
    \item The spending (or expenditure) multiplier
    \item The investment multiplier
    \item The balanced budget multiplier
    \item The aggregate demand multiplier
\end{enumerate}
\item Which of the following could cause the aggregate demand curve to shift to the left?
\begin{enumerate}[label=(\alph*)]
    \item An increase in government spending.
    \item A decrease in interest rates.
    \item Contractionary demand management policies.
    \item An increase in consumer confidence.
    \item An increase in the money supply.
\end{enumerate}
\item In 2004 the United States had a trade deficit of \$603 billion; therefore
\begin{enumerate}[label=(\alph*)]
    \item The value of the dollar increased significantly.
    \item The United States' GDP decreased by \$603 billion.
    \item The United States imported exactly \$603 billion worth of goods.
    \item The United States had a trade surplus.
    \item Americans consumed more than they produced.
\end{enumerate}
\item Which of the following is true?
\begin{enumerate}[label=(\alph*)]
    \item Average total cost = total fixed costs plus total variable costs.
    \item Average total cost = total costs divided by the marginal cost of production.
    \item Average total cost = total variable costs divided by the number of units produced.
    \item Average total cost = total fixed costs divided by the number of units produced.
    \item Average total cost = average variable cost plus marginal cost.
\end{enumerate}
\item For a given level of government spending the federal government usually experiences a budget\_\_\_\_during economic\_\_\_\_and a budget \_\_\_\_\_\_during economic\_\_\_\_\_\_\_.
\begin{enumerate}[label=(\alph*)]
    \item surplus expansion surplus expansion
    \item deficit recession deficit expansion
    \item surplus expansion deficit recession
    \item deficit expansion deficit recession
    \item surplus expansion surplus recession
\end{enumerate}
\end{enumerate}

\section*{True / False (10 marks)}
\textit{Answer True or False}
\begin{enumerate}[label=\arabic*.]
\setcounter{enumi}{5}
\item True or False: The a and c coefficients must be significant for bi-directional feedback to be established in a VAR model.
\item True or False: A stronger stock market typically causes a decrease in the consumption function while increasing aggregate demand.
\item True or False: The spending multiplier effect is enhanced when the price level remains constant during a decrease in aggregate demand.
\item True or False: A decrease in the price of a particular product will lead to a decrease in the quantity demanded by consumers.
\item True or False: Investment is always greater than savings in a private closed economy.
\end{enumerate}

\section*{Long Form Response (20 marks)}
\textit{Answer each question with a clear and complete explanation.}
\begin{enumerate}[label=\arabic*.]
\setcounter{enumi}{10}
\item Discuss the implications of a bank acquiring an additional \$1 in deposits without a required reserve ratio, and analyze how this could lead to an infinite expansion of the money supply.
\item Discuss the impact of trade between an industrial country abundant in capital and a rural country abundant in land on the prices of these resources, considering comparative advantage.
\end{enumerate}

\vfill
\noindent\textit{End of Paper}
\end{document}